\begin{abstract}
\normalsize
Retrieval-Augmented Generation (RAG) has become the major method for grounding Large Language Model (LLM) outputs with external evidence,
but existing text-only pipelines cannot accommodate multimodal information.
Multimodal RAG systems are being developed for use with Multimodal LLMs (MLLMs), and can process texts and images.
Current systems flatten visual features into text captions, or rely on CLIP-based embeddings that prioritize semantic alignment
over fine-grained visual details, limiting RAG in applications that require visual precision.

This project presents a multimodal RAG pipeline that preserves fine-grained visual features while enabling cross-modal alignment.
We propose a dual-representation framework maintaining separate representation spaces, 
using state-of-the-art pretrained foundational models for feature extraction and cross-modal alignment.
Our contributions include:
(i) a complete multimodal RAG pipeline that allows for text-only and text-image queries, with retrieved evidence covering both modalities;
(ii) a method without fine-tuning requirements that leverages state-of-the-art foundational models to use in evidence retrieval;
and (iii) empirical evaluation on the effect of different retrieval modes in grounding generation to retrieved evidence.
We compare dense retrieval against graph-based retrieval and analyze how different retrieval modes affect grounding in MLLM generation.
\end{abstract}